\begin{lemma}[Reflexivity of the lifted relation]
If $r$ is a preorder on $Q$, then
\[
\forall X\in\mathsf{PowerQ}(Q),\quad
\mathsf{PowerRel}(r)(X,X).
\].
\end{lemma}
\begin{proof}
For each $x\in\mathsf{PowerQ}(Q)$, let $P(x)$ be the proposition
\(\mathsf{PowerRel}(r)(x,x)\).  We prove $P(x)$ by well-founded induction
on the child-to-parent relation.  This induction is legitimate by the
well-foundedness statement \noderef{PowerChildWellFounded}; thus, in the
induction case for $x$, we may use $P(y)$ whenever
$\mathsf{PowerChild}(y,x)$.

First suppose that $x=\mathsf{atom}(q)$.  The atom--atom reduction
\noderef{PowerRelAtomAtom}, obtained by unfolding the recursive definition
of $\mathsf{PowerRel}$ at the two atoms, identifies the goal
\(\mathsf{PowerRel}(r)(\mathsf{atom}(q),\mathsf{atom}(q))\) with the goal
$r(q,q)$.  The stated instance $[\mathsf{IsPreorder}\ Q\ r]$ supplies
reflexivity, so $r(q,q)$ holds and hence $P(\mathsf{atom}(q))$ holds.

Now suppose that $x=\mathsf{node}(\iota,h,f)$.  By the node--node reduction
\noderef{PowerRelNodeNode}, after unfolding the recursive definition, it is
enough to prove
\[
\forall i\in\iota\;\exists j\in\iota,
\quad \mathsf{PowerRel}(r)(f(i),f(j)).
\]
Fix an arbitrary $i\in\iota$.  Choose the right-hand index to be the same
index, namely $j=i$.  By the child definition \noderef{PowerChild}, the
term
\[
\langle i,\mathord{rfl}\rangle
\]
is a witness of
$\mathsf{PowerChild}(f(i),\mathsf{node}(\iota,h,f))$: it supplies the
index $i$ and the required equality $f(i)=f(i)$.  The induction hypothesis
therefore gives
\[
\mathsf{PowerRel}(r)(f(i),f(i)).
\]
Combining this proof with the chosen witness $j=i$ gives the existential
claim for the fixed $i$.  Since $i$ was arbitrary, universal introduction
gives the displayed $\forall i\,\exists j$ statement, and the reduction
then gives $P(\mathsf{node}(\iota,h,f))$.
\end{proof}
